% id: 86-26
% book: 베이직쎈 미적분1 · 05 도함수의 활용 (2)
% section: 개념 39 함수의 최대와 최소의 활용
% figure: tikz:fig-86-26
% answer: 
% answer_source: 
% variant_level: 0 (원본 전사)
% 이 파일은 build.mjs 가 items.json 에서 만든다. 고칠 때는 items.json 을 고친다.
\smash{\raisebox{1.5mm}{\dmkichul{베이직쎈 미적분1 86쪽 26번}}}
\noindent\begin{minipage}[t]{0.58\linewidth}\vspace{0pt}\raggedright 오른쪽 그림과 같이 함수 $y=3-x^2$의 그래프와 $x$축으로 둘러싸인 도형에 내접하고, 한 변이 $x$축 위에 있는 직사각형~$\pt{ABCD}$가 있다. 제1사분면 위의 점~$\pt{A}$의 $x$좌표를 $a$라 할 때, 다음에 답하시오.\end{minipage}\hfill\begin{minipage}[t]{0.4\linewidth}\vspace{0pt}\centering \resizebox{\ifdim\width>\linewidth\linewidth\else\width\fi}{!}{% 86-26(86쪽) — 포물선 $y=3-x^2$ 과 $x$축으로 둘러싸인 도형에 내접한 직사각형 ABCD. 스캔 화질이 낮아 TikZ 로 다시 그림.
% 원문에 있는 것만: 라벨 3, y=3-x^2, A, B, C, D, O, x, y 와 회색으로 칠한 직사각형. a 의 값은 원문에 없으므로 적지 않는다.
\begin{tikzpicture}[x=0.48cm, y=0.38cm, line width=0.5pt]
  \fill[black!15] (-1,0) rectangle (1,2);
  \draw[->, >=stealth, line width=0.7pt] (-2.5,0) -- (2.95,0) node[below=1pt, font=\small] {$x$};
  \draw[->, >=stealth, line width=0.7pt] (0,-1.6) -- (0,3.95) node[left=1pt, font=\small] {$y$};
  \draw[line width=0.6pt, smooth] plot coordinates {(-2.05,-1.2) (-1.8,-0.24) (-1.5,0.75) (-1.2,1.56) (-0.9,2.19) (-0.6,2.64) (-0.3,2.91) (0,3) (0.3,2.91) (0.6,2.64) (0.9,2.19) (1.2,1.56) (1.5,0.75) (1.8,-0.24) (2.05,-1.2)};
  \draw[line width=0.5pt] (-1,0) rectangle (1,2);
  \node[left=1pt, font=\small] at (0,3) {$3$};
  \node[anchor=west, font=\small] at (0.35,3.05) {$y=3-x^2$};
  \node[left=2pt, font=\small] at (-1,2) {$\pt{B}$};
  \node[right=2pt, font=\small] at (1,2) {$\pt{A}$};
  \node[below=1pt, font=\small] at (-1,0) {$\pt{C}$};
  \node[below left=-2pt and -3pt, font=\small] at (0,0) {$\pt{O}$};
  \node[below=1pt, font=\small] at (1,0) {$\pt{D}$};
\end{tikzpicture}
}\end{minipage}\par
\dmsubprob{1} $\seg{AD}$, $\seg{CD}$의 길이를 $a$에 대한 식으로 나타내시오.
\dmsubprob{2} $a$의 값의 범위를 구하시오.
\dmsubprob{3} 직사각형~$\pt{ABCD}$의 넓이를 $S(a)$라 할 때, $S(a)$를 구하시오.
\dmsubprob{4} 직사각형~$\pt{ABCD}$의 넓이의 최댓값을 구하시오.
